\subsection{Diagonalisable endomorphisms}

DEFINITION 5. --- Let E be a finite dimensional vector space over a commutative field K and let F be a set of endomorphisms of E. Then F is said to be diagonal with respect to a basis \((e_{i})\) of E if the matrix of each \(u \in F\) with respect to \((e_{i})\) is diagonal. The set F is said to be diagonalisable if there exists a basis of E with respect to which F is diagonal.

This definition applies in particular to the case when \(\mathfrak{F}\) contains only one element \(u\) ; we then say that \(u\) is diagonal (diagonalisable). Note also that \(\mathfrak{F}\) is diagonal with respect to a basis \((e_i)\) if and only if the \((e_i)\) are common eigenvectors of all the elements of \(\mathfrak{F}\) ; it follows that \(\mathfrak{F}\) is diagonalisable if and only if \(E\) is generated by eigenvectors common to all the elements of \(\mathfrak{F}\) .

Let A be a subalgebra of \(\operatorname{End}_{\mathbf{K}}(\mathbf{E})\) containing \(Id_{E}\) . Then A is diagonalisable if and only if it is isomorphic to an algebra \(K^{r}\) (in other words is diagonalisable in the sense of V, p.~28, Def. 1); indeed, if A is isomorphic to \(K^{r}\) , then A is diagonalisable by V, p.~29, Prop. 1; conversely, if A is diagonalisable, then it is isomorphic to a subalgebra of the algebra of diagonal matrices, which is isomorphic as an algebra to \(K^{n}\) , \(n = \dim(\mathbf{E})\) , hence A is isomorphic to some algebra \(K^{r}\) (V, p.~30, Prop. 3).

PROPOSITION 12. --- Let E be a finite dimensional vector space over a commutative field K, and let u be an endomorphism of E. Then the following conditions are equivalent :

\begin{enumerate}
\def\labelenumi{(\roman{enumi})}
\item
  \(u\) is diagonalisable.
\item
  E is the direct sum of the eigenspaces of \(u\) .
\item
  All the roots of the minimal polynomial of \(u\) are in \(K\) , and these roots are all simple.
\end{enumerate}

Moreover, if these conditions are satisfied, then every subspace of E closed under u is the direct sum of its intersections with the eigenspaces of u.

The equivalence of (i) and (ii) follows from the preceding remarks and VII, p.~31, Prop. 2. Suppose u is diagonalisable, and let \((\alpha_{i})\) be its family of eigenvalues, and \((\mathbf{V}_{i})\) the corresponding family of eigenspaces; since the restriction of u to \(V_{i}\) is the homothety defined by \(\alpha_{i}\) , it annihilates the polynomial \(X - \alpha_{i}\) ; it follows that u annihilates the polynomial \(\prod_{i}(X - \alpha_{i})\) , which is therefore a multiple of the minimal polynomial of u, and hence coincides with it, which proves (iii). Conversely, if (iii) is satisfied then there exists a basis of E with respect to which the matrix U of u is a diagonal block of Jordan matrices \(U_{m,\alpha}\) (VII, p.~35, Prop. 8); then by Prop. 9 the integers m are all equal to 1 and so U is diagonal. Finally, the last assertion follows from the fact that if u is diagonalisable then its eigenspaces are the primary components of \(E_{u}\) , and from VII, p.~9, Cor. 1.

COROLLARY. --- If all the roots of the characteristic polynomial of u are in K, and they are all simple, then u is diagonalisable.

Indeed the minimal polynomial divides the characteristic polynomial.

PROPOSITION 13. --- Let E be a finite dimensional vector space over a commutative field K, let F be a set of endomorphisms of E, and let A be the subalgebra of \(\mathrm{End}_{\mathbf{K}}(\mathbf{E})\) generated by F and \(\mathrm{Id}_{\mathbf{E}}\) . Then the following conditions are equivalent:

\begin{enumerate}
\def\labelenumi{(\roman{enumi})}
\item
  \(\mathfrak{F}\) is diagonalisable.
\item
  The K-algebra \(\mathbf{A}\) is diagonalisable.
\item
  The elements of \(\mathfrak{F}\) are diagonalisable and commute with one another.
\end{enumerate}

If \((e_{i})\) is a basis of E with respect to which F is diagonal, then A is contained in the algebra of endomorphisms which are diagonal with respect to this basis, so is also diagonalisable; if A is diagonalisable, then the same argument shows that F is diagonalisable. This shows the equivalence of (i) and (ii). Since any two diagonal matrices commute, we have \((\mathrm{i}) \Rightarrow (\mathrm{iii})\) , and it remains to prove the converse. Suppose then that the elements of F are diagonalisable and commute with each other. We will make use of the following lemma:

Lemma 3. --- Let g and h be two commuting endomorphisms of a vector space E. Then each eigenspace of g is closed under h.

Indeed, if \(\mathbf{W}_{\lambda}\) is the eigenspace of \(g\) corresponding to the eigenvalue \(\lambda\) , then for all \(x \in \mathbf{W}_{\lambda}\) we have

\[
g h. x = h g. x = h. \lambda x = \lambda h. x,
\]

which says that \(h \cdot x \in \mathbf{W}_{\lambda}\) .

Let us now return to the proof of Prop. 13. Among all decompositions of E as a direct sum of nonzero subspaces each closed under all the elements of \(\mathfrak{F}\) , choose one with the greatest number of components (the dimension of E is an upper bound for this number), say \(E = \sum_{i\in I}E_i\) . Let \(u\in \mathfrak{F}\) and let \(E = \sum_{\alpha}V_{\alpha}\) be the

decomposition of E as the direct sum of the eigenspaces of u. By Lemma 3, each \(V_{\alpha}\) is closed under F, and hence so is each \(V_{\alpha} \cap E_{i}\) ; by Prop. 12 each \(E_{i}\) is the direct sum of the \(V_{\alpha} \cap E_{i}\) . The choice of the \(E_{i}\) thus forces each \(E_{i}\) to be contained in one of the \(V_{\alpha}\) ; thus the restriction of u to each \(E_{i}\) is a homothety. Since this is true for all the elements of F, it follows that F is diagonalisable.

COROLLARY. --- The sum and the composite of two commuting diagonalisable endomorphisms of E are diagonalisable.

